\chapter*{Abstract}

The identification of tumors in tissues is currently based on histopathological examination, during which biopsy samples are prepared, treated with histological stains and evaluated by experienced pathologists. Although this approach represents the valid clinical reference, it requires specific sample preparation, specialized skills and extended analysis times. In addition, the evaluation can include a subjective/expert-dependent influence, which plays a role especially when the samples are particularly heterogeneous, inflamed or have reactive stroma. The development of quantitative and reproducible methods, able to support and, in the future, partially automate histological inspection, represents an important direction for biomedical research. 

Raman spectroscopy emerges as a promising technique in this scenario, as it provides chemically specific information, does not require the use of markers and is non-destructive. The high information content of Raman spectra, due to the complexity of biological samples, makes it necessary to resort to computational and machine learning tools suitable for interpreting and classifying complex, high-dimensional data. 

This thesis develops a machine learning pipeline for the analysis of spontaneous Raman hyperspectral maps acquired from patient-delivered head and neck squamous cell carcinoma (HNSCC) biopsies. The work combines Principal Component Analysis (PCA) and the training of an Artificial Neural Network (ANN), respectively, for dimensionality reduction and supervised classification of Raman data, aimed at identifying tumor tissue. Overall, the results highlight the potential of Raman spectroscopy in capturing biochemical differences and show how machine learning can support the development of computational procedures for the characterization of tumor tissues.

\vspace{1em}
\noindent\textbf{Keywords:} Raman spectroscopy, hyperspectral maps, machine learning, tumor tissue characterization.
